\section{A capacity comparison on one period}\label{sec:capacity}

The probability of returning to a pin includes arbitrarily many turns
of the fresh walk.  We compare this probability directly, through a
hitting problem on one cut.  The local conclusion needed for
magnetization is the following uniform comparison.  Write
$B_r(x)=\{z\in\Z^3:|z-x|_\infty\le r\}$ for a site ball and set
$R_*=\lceil\beta^{5/3}\rceil$.

\begin{proposition}[Uniform comparison away from pins]\label{prop:comparison}
For every fixed positive half-integer $S$ there are
$\beta_0(S)<\infty$ and a function $\varepsilon_S(\beta)\to0$
as $\beta\to\infty$ such that the following holds for every
$\beta\ge\beta_0(S)$.
Let $\Lambda$ be any finite cube, let $E=D_0\cup F$ with $F$ any
finite deterministic set of slot-time pins, and let $i=(x,a)$ be a
time-zero post-seam slot.  Suppose that $B_{50R_*}(x)$ is interior
to $\Lambda$ and that $E$ has no point above this ball.  Then
\begin{equation}\label{eq:pinnedcomparison}
 \abs{d_E(i)-\frac1\ell n_S^{\rm sw}(\beta)}
 \le \varepsilon_S(\beta)\beta^{-5/2}.
\end{equation}
The same uniform error bound, after enlarging $\varepsilon_S$ if
necessary, holds for
$\abs{\EE_{E\cup\{i\}}q_i-\ell^{-1}n_S^{\rm sw}(\beta)}$.
\end{proposition}

The bound is uniform over the number and locations of all pins
outside the displayed ball.  This uniformity permits a field to be
represented by additional random pins.  The proof is completed in
Section~\ref{sec:pressure}, after the required residual estimates
have been established.

We first give a finite-dimensional criterion that identifies the two
estimates needed for this comparison: a scalar residual at a trial
potential and a residual bound in the energy seminorm.  Both directions
of the period occur because its transition matrix need not be symmetric.
The relation between hitting potentials and energy also appears in
nonreversible potential theory; see
\cite[Section~2, Equation~(2.6)]{GaudilliereLandim2014}.
The comparison below follows from direct finite-dimensional identities.

Let $H$ be a finite coordinate set, let $K\subset H$ be killing
coordinates, and fix $i\in H\setminus K$.  Suppose that $U$ is a
nonnegative doubly substochastic matrix on $\R^H$, with zero rows and
columns on $K$.  It acts forward on columns.  Complete each column to
one by adding a cemetery state, and use $U$ for the successive
transitions of the resulting killed chain.  Let $q$ be its probability,
starting at $i$, of returning to $i$ after at least one step before
killing.  Let $f_j$ be the probability of hitting $i$ before killing
from $j$, with $f_i=1$ and $f=0$ on $K$.  Thus $0\le f\le1$.

Write
\[
 A=I-U,\qquad \eta=\ind_{\{i\}},\qquad
 e_z=\ip{z}{Az}.
\]
All inner products and norms use counting measure.  Since $U$ is an
$\ell^2$ contraction, $e_z\ge0$, and real inner products give
$e_z=\ip{z}{A^*z}$.  In the application, $H$ is the set of holes at
the time-zero cut and $U$ is one traversal killed on $E$.
Lemma~\ref{lem:kernelenergy} supplies exactly these contraction
properties and their continuous-time energy interpretation.

Choose a real vector $h$ with $h_i=1$ and $h=0$ on $K$, and a real
number $c$.  Define its two residuals by
\begin{equation}\label{eq:capacityresiduals}
 \mu_+=Uh-h+c\eta,\qquad
 \mu_-=U^*h-h+c\eta.
\end{equation}
Thus $Ah=c\eta-\mu_+$ and $A^*h=c\eta-\mu_-$.

\begin{lemma}[Finite capacity comparison]\label{lem:capacity}
For the preceding data, put $y=f-h$.  Then
\begin{align}
 A^*f&=(1-q)\eta,                                      \label{eq:harmonic}\\
 e_y&=\ip{y}{\mu_-},                                  \label{eq:errorenergy}\\
 q-(1-c)
 &=\ip{h}{\mu_+}+\ip{y}{\mu_+}.                     \label{eq:capacityidentity}
\end{align}
No assumption of almost-sure killing is required.

Suppose, in addition, that $\norm{h}_\infty\le M$ and there are
numbers $X,\varepsilon\ge0$ such that, for both $d=+,-$,
\begin{equation}\label{eq:abstractdual}
 \abs{\ip{\mu_d}{z}}\le X\sqrt{e_z}+\varepsilon
 \quad\text{whenever }z|_K=0\text{ and }\norm{z}_\infty\le1+M.
\end{equation}
Then
\begin{equation}\label{eq:capacitystability}
 \abs{q-(1-c)}
 \le \abs{\ip{h}{\mu_+}}+2X^2+2\varepsilon.
\end{equation}
\end{lemma}

\begin{proof}
For $j\notin K\cup\{i\}$, the first-step decomposition of the
hitting event gives $f_j=\sum_kU_{kj}f_k=(U^*f)_j$.
The same equation is zero on $K$, whereas $(U^*f)_i=q$.
This proves \eqref{eq:harmonic}, even if the chain can remain forever
in a class disjoint from $i$ and the cemetery state.
In particular,
\[
 \ip{f}{Af}=\ip{A^*f}{f}=1-q.
\]
Subtracting the adjoint residual equation from
\eqref{eq:harmonic} gives
\[
 A^*y=(1-q-c)\eta+\mu_-.
\]
Since $y_i=0$, pairing with $y$ proves \eqref{eq:errorenergy}.
To obtain the capacity identity, use $h_i=f_i=1$ to write
\[
 1-q=\ip{A^*f}{h}=\ip{f}{Ah}
     =c-\ip{f}{\mu_+}.
\]
Substituting $f=h+y$ proves \eqref{eq:capacityidentity}.

The vector $y$ is admissible in \eqref{eq:abstractdual}.
Hence $e_y\le X\sqrt{e_y}+\varepsilon$, which implies
$\sqrt{e_y}\le X+\sqrt\varepsilon$.  Consequently
\[
 \abs{\ip{y}{\mu_+}}
 \le X^2+X\sqrt\varepsilon+\varepsilon
 \le2X^2+2\varepsilon.
\]
Insert this bound into \eqref{eq:capacityidentity}.
\end{proof}

The adjoint residual in \eqref{eq:errorenergy} controls the distance
between the hitting potential and the trial vector.  The forward
residual then controls the capacity itself.  In particular, one scalar
estimate and two energy estimates suffice; no expansion in the number
of returns is involved.

The same criterion can be averaged over a random environment.
Precisely, suppose that all preceding objects are measurable, that
$c\in[0,1]$ is deterministic, and that \eqref{eq:abstractdual} holds
on a measurable event $\mathcal G$ with a uniform bound $M$.
Then \eqref{eq:capacitystability} and $0\le q\le1$ give
\begin{equation}\label{eq:averagedcapacity}
 \abs{\EE q-(1-c)}
 \le \EE\!\left[\ind_{\mathcal G}\abs{\ip{h}{\mu_+}}\right]
      +2\EE[\ind_{\mathcal G}X^2]
      +2\EE[\ind_{\mathcal G}\varepsilon]
      +\PP(\mathcal G^c).
\end{equation}
We next construct an event on which diffusion and spatial flow
estimates are available.  After that construction, the all-turn free
potential will provide the trial vector $h$.
